\section{A closed coefficient-representation chain by Gaussian elimination}
\label{sec:highest-weight-closure}

\status{HP1--HP6 remove the retained characteristic-zero polynomial
highest-weight inputs from MR5/MR19 and AR13--AR14.  The proof applies to
the literal polynomial module used in the source, over every field of
characteristic zero.  It proves the two highest-fixed-vector interfaces,
absolute tensor irreducibility, classification, and the precise
contragredient and scalar-extension statements, then reaches the full
\texttt{wt\_weights} endpoint on its original data.  The argument uses
finite polynomial equations, Gaussian elimination, and a factorial
pairing.  It does not assume a Schur-module irreducibility theorem,
highest-weight classification, or a noncompact representation theorem.
The compatible geometric/theta/local state and the separate automorphic
Galois inputs remain the frontier of the larger form construction.}

We use indices $1,\ldots,n$ in this section.  In the source they are
$0,\ldots,n-1$, with the evident shift.  If $\lambda_1\ge\cdots\ge\lambda_n$
are integers, put $d=\lambda_n$ and $\nu_i=\lambda_i-d$; for $n=0$ use
$d=0$.  Let $X=(x_{ij})$ be the matrix of polynomial variables.  Write
$U_+$ and $U_-$ for the upper and lower unitriangular groups and $T$ for
the diagonal group.  Define
\[
 \mathcal N_\nu(K)=\{f\in K[x_{ij}]:
 f(Xt)=t_1^{\nu_1}\cdots t_n^{\nu_n}f(X),\quad
 f(Xu)=f(X)\ (u\in U_+)\}.
 \tag{HP.1}
\]
The first equality means column multihomogeneity.  Both equalities can
be read as polynomial identities.  Over an infinite field, equality at
all indicated points is equivalent to this reading.

This is exactly the source's \texttt{dualWeyl K n n nu}, hence
\texttt{nablaSq K n lambda} for $\nu=\lambda.\mathrm{natW}$.
The action, including its transpose and determinant, is
\[
       (\rho_\lambda(g)f)(X)=(\det g)^d f(g^tX).
 \tag{HP.2}
\]
It is \texttt{sqAct}, not substitution by $gX$ or $g^{-1}X$.
For a finite index set $I$ and weights $\lambda_\tau$, use separate
matrix variables $X_\tau$ and take the image of the tensor product
$\bigotimes_{\tau\in I}\mathcal N_{\nu_\tau}(K)$ under multiplication.
The polynomial monomial bases in disjoint variables show that this
map is injective.  Its image is the literal \texttt{VV K lambda};
(HP.2) factor by factor is its \texttt{rhoK} action.  Empty $I$ or
$n=0$ gives $K$, with highest vector $1$.

\paragraph{HP1: the literal fixed line and its rational structure.}
\label{claim:HP1}
Every $\mathcal N_\nu(K)$ is the solution of a finite system of linear
equations with integer coefficients inside a finite polynomial space.
Indeed, column multidegree $\nu$ forces total degree $N=\sum\nu_i$.
For each $i<j$, substitute column $j+t$ column $i$ for column $j$ and
set all coefficients of $t^a x^\alpha$ in the difference to zero.
There are finitely many equations because its degree is at most $N$.
The elementary column operations generate $U_+$.  These are the same
equations as the generic \texttt{colAddGen} kernels in (HP.1).
Thus, in characteristic zero, row reduction over $\mathbb Q$ gives
\[
 K\otimes_{\mathbb Q}\mathcal N_\nu(\mathbb Q)
       \simeq\mathcal N_\nu(K).
 \tag{HP.3}
\]
No integral freeness or equality between the integral Weyl and dual
Weyl lattices is needed or asserted.  The same argument applies to
the tensor polynomial space; it also identifies its scalar extension
with the same polynomial subspace over the new field.  Moreover, the
joint equations in separate blocks describe exactly that tensor space:
expand in a basis of the other blocks, apply one block's equations to
each coefficient, and repeat.  Equivalently, choose a complement to
each one-block kernel in its finite polynomial space; the tensor
decomposition then shows that a term with any complementary factor
cannot satisfy all the block equations.

Let $\Delta_r(X)$ be the leading $r$ by $r$ minor and set
$\nu_{n+1}=0$.  Then
\[
             h_\nu(X)=\prod_{r=1}^n
                   \Delta_r(X)^{\nu_r-\nu_{r+1}}
 \tag{HP.4}
\]
belongs to $\mathcal N_\nu(K)$, is nonzero, and is fixed by
$\rho_\lambda(U_+)$.  Each principal minor is unchanged by lower
unitriangular row operations and by upper unitriangular column
operations.  The column weights telescope to $\nu$, and
$h_\nu(1)=1$.  The determinant twist is $1$ on $U_+$.
The diagonal action on this vector is the character $\lambda$:
\[
 \rho_\lambda(\operatorname{diag}t)h_\nu
                      =\prod_i t_i^{\lambda_i}h_\nu.
 \tag{HP.5}
\]

Conversely, suppose $f\in\mathcal N_\nu(K)$ is fixed by
$\rho_\lambda(U_+)$.  This says $f(lX)=f(X)$ for $l\in U_-$.
On the open set where every $\Delta_r(X)$ is nonzero, Gaussian
elimination gives $X=l a u$ with $l\in U_-$, $a$ diagonal and
$u\in U_+$, and $a_i=\Delta_i/\Delta_{i-1}$, where $\Delta_0=1$.
Consequently
\[
 f(X)=f(a)=f(1)\prod_i a_i^{\nu_i}
                =f(1)h_\nu(X).                    \tag{HP.6}
\]
This open set is dense over any infinite field: a polynomial vanishing
there has product with a suitable power of $\prod_r\Delta_r$ equal
to zero, and the polynomial ring is a domain.  Equivalently, its
Gaussian coordinates give an open rational parametrization and the
usual polynomial identity principle.  Hence (HP.6) holds identically.
For several blocks apply Gaussian elimination in every block.  Thus
\[
 (\mathrm{VV}_K(\lambda))^{\prod U_+}
       =K h,\qquad h=\prod_{\tau\in I}h_{\nu_\tau}(X_\tau).
 \tag{HP.7}
\]
This is the exact source \texttt{HwFixedLine}; its distinguished vector
is \texttt{mtwHw}, not a highest vector chosen after an isomorphism.

We will also use the same argument before dominance is known.
If $\nu\in\mathbb N^n$ is arbitrary and a nonzero
$\mathcal N_\nu(K)$ contains a nonzero left-$U_-$-fixed polynomial,
(HP.6) holds as an equality of rational functions and forces
$\nu_r\ge\nu_{r+1}$.  For $r<n$, take the identity matrix except
for the rows/columns $r,r+1$, where the block is
$\left(\begin{smallmatrix}t&-1\\1&0\end{smallmatrix}\right)$.
All its leading minors are $1$, except $\Delta_r=t$.
Then (HP.6) reads $f(X_t)=f(1)t^{\nu_r-\nu_{r+1}}$ for $t\ne0$.
A polynomial cannot equal a negative power times a nonzero constant.
Moreover $f(1)\ne0$ by the identity principle.  This proves the claimed
dominance without invoking a classification theorem.

\paragraph{HP2: fixed-vector existence and absolute irreducibility.}
\label{claim:HP2}
Let $S$ be a nonzero subspace of $\mathrm{VV}_K(\lambda)$ stable under
$G=\prod_{\tau\in I}\mathrm{GL}_n(K)$, with $\operatorname{char}K=0$.
The diagonal action has finitely many weights.  For the polynomial
model these are the row multidegrees, shifted by $d_\tau$ in every
coordinate of block $\tau$.  The weight projections preserve $S$:
choose a diagonal element at which these finitely many different
Laurent characters have distinct values, and interpolate its spectral
projections by polynomials.  Such an element exists because a finite
product of nonzero Laurent polynomials cannot vanish on the whole
torus over an infinite field.

Choose a nonzero vector $v\in S$ of weight $\mu$ maximizing
$H(\mu)=\sum_{\tau,i}(n-i)\mu_{\tau,i}$ among weights occurring in $S$.
For an upper elementary matrix $g=1+tE_{ij}$, $i<j$, the polynomial
$\rho_\lambda(g)v$ has coefficients $v_a$ of weights
$\mu+a(e_{\tau,i}-e_{\tau,j})$.  This follows directly from replacing
row $j$ of $X_\tau$ by row $j+t$ row $i$; the determinant twist is $1$.
All $v_a$ belong to $S$, since coefficients are linear combinations of
values at sufficiently many distinct $t\in K$.  For $a>0$ their heights
are strictly greater than $H(\mu)$, so they vanish.  Hence $v$ is fixed
by every elementary upper unipotent and therefore by $\prod U_+$.
Together with HP1 this proves
\[
             S\ne0\quad\Longrightarrow\quad h\in S.
 \tag{HP.8}
\]
In particular this proves the exact \texttt{HwFixedEx} interface,
even over any characteristic-zero field.  The proof also works for
$\mathcal N_\nu$ with arbitrary nonnegative $\nu$; with HP1 it shows
that a nonzero such space necessarily has dominant $\nu$.

To show that (HP.8) implies simplicity here, the required duality
must be justified.  On the polynomial ring define the bilinear
factorial pairing
\[
        B(f,g)=\sum_\alpha \alpha! f_\alpha g_\alpha,
       \qquad \alpha!=\prod_j\alpha_j! .             \tag{HP.9}
\]
It has no complex conjugation.  Its restriction to
$\mathrm{VV}_{\mathbb Q}(\lambda)$ is nondegenerate: every nonzero
rational polynomial $f$ has $B(f,f)>0$ when read in $\mathbb R$.
A rational basis therefore has a positive-definite rational Gram
matrix, whose nonzero rational determinant remains nonzero in every
characteristic-zero field.  Equation (HP.3) proves nondegeneracy on
the entire $\mathrm{VV}_K(\lambda)$, not merely on an arbitrary
subspace of polynomials.  In particular $B(h,h)$ is a strictly
positive integer and remains nonzero in $K$.

The substitution identity for this pairing is
\[
              B(\rho_\lambda(g)f,z)
                       =B(f,\rho_\lambda(g^t)z).     \tag{HP.10}
\]
One direct verification writes
$B(f,z)=f(\partial)z|_{X=0}$ and uses the linear chain rule; checking
monomials gives the same multinomial coefficients $\alpha!$.
Both determinant factors are $(\det g)^{d_\tau}$, so the identity
holds for negative determinant twists as well as nonnegative ones.
It holds block by block for the tensor product.

If $S$ is proper, its annihilator $S^{\perp_B}$ is nonzero by
nondegeneracy.  It is stable under $G$, because for $x\in S^{\perp_B}$
and $y\in S$,
$B(\rho(g)x,y)=B(x,\rho(g^t)y)=0$.
Equation (HP.8) now puts $h$ in both $S$ and $S^{\perp_B}$,
contradicting $B(h,h)\ne0$.  Thus $S$ is the whole space.  This proves
simplicity over every characteristic-zero field and, by (HP.3), after
every extension of that field.  It closes the dual-Weyl, tensor and
highest-fixed-line chain with finite-dimensional polynomial arguments.

\paragraph{HP3: the actual rational unitary group and every coefficient field.}
\label{claim:HP3}
Let $E$ be a CM field, $c$ its specified involution, and let
$\sigma_\tau:E\hookrightarrow K$ be pairwise distinct embeddings
with $\sigma_\tau\ne\sigma_{\tau'}c$ for all $\tau,\tau'$.
Then the representation
\[
 \gamma\longmapsto
   \rho_K((\sigma_\tau(\gamma))_\tau)
       \quad(\gamma\in U_n(E/E^+)(\mathbb Q))          \tag{HP.11}
\]
on the same $\mathrm{VV}_K(\lambda)$ is absolutely irreducible.
Here the rational notation denotes the rational points of the
restriction of scalars from $E^+$ to $\mathbb Q$, exactly the source's
\texttt{GQ}, with its injective finite-adelic realization forgotten.

For completeness the density assertion needed in this passage has a
short proof.  The rational vector space
$\mathfrak s=\{X\in M_n(E):X+{}^t c(X)=0\}$ has the Cayley map
\[
                 X\longmapsto(1-X)(1+X)^{-1}.        \tag{HP.12}
\]
It is a rational isomorphism between nonempty open subsets of
$\mathfrak s$ and the rational unitary group, with the same inverse
formula.  In rational coordinates a polynomial over any extension
field which vanishes on all $\mathbb Q$-points is zero: induct on the
number of coordinates and use infinitely many rational values in one
coordinate.  Removing the finitely many Cayley denominators preserves
this argument.  Hence the rational points are Zariski dense, including
after extension to an algebraic closure $\overline K$.

Over $\overline K$, the pairs of conjugate embeddings of $E$ give
\[
 (\operatorname{Res}_{E^+/\mathbb Q}U_n)_{\overline K}
                \simeq\prod_{\{\sigma,\sigma c\}}\mathrm{GL}_n.
 \tag{HP.13}
\]
Indeed the unitary equation reads $A_{\sigma c}=A_\sigma^{-t}$,
so one matrix in each pair is arbitrary.  Projection to the selected
$\sigma_\tau$ is surjective by the distinct-nonconjugate hypothesis.
A stable subspace is defined by the equations obtained by composing
its inclusion, the action matrix, and its quotient map.  These are
regular equations in the group entries; density makes them hold on
all of (HP.13).  Its scalar extension is consequently stable under
all selected general linear groups and is zero or full by HP2.
Faithfulness of scalar extension gives the same conclusion over $K$.
The argument applies after every $K\hookrightarrow K'$, since
injectivity preserves the distinctness conditions.  It proves the
literal \texttt{tenRepF\_isIrreducible} and its base-change clauses;
no embedding of all of $K$ into $\mathbb C$ is assumed.

The action matrices are algebraic in exactly the source's sense.
The polynomial substitution part is polynomial in matrix entries.
A negative power of $\det\sigma_\tau(\gamma)$ is replaced, on the
unitary group, by the corresponding nonnegative power of
$\det\sigma_\tau c(\gamma)$.  Express every entry of $\gamma$ in a
fixed $\mathbb Q$-basis of $E$; each embedding is a linear polynomial
in those coordinates.  This gives the finite polynomial matrix
presentation in \texttt{IsAlgWeightF}.  In the dual representation
replace $\gamma^{-1}$ by ${}^t c(\gamma)$ and transpose its action
matrix.  Its coefficients are again polynomials in the same rational
coordinates, over the original $K$.

\paragraph{HP4: classification and the source's dual conventions.}
\label{claim:HP4}
We first derive the classification used in AR13, rather than leaving
it as an additional representation theorem.  An algebraic matrix
representation of the rational unitary points extends to the split
product (HP.13): its identity and multiplication equations hold on
dense rational points, hence identically.  Its matrix entries are
regular functions of the general linear matrices, thus polynomials
in their entries and inverse determinants.

For an arbitrary finite-dimensional rational representation of a
product of general linear groups over $\mathbb C$, diagonal torus
restriction has a finite direct sum of integral weights.  This can
be seen without simultaneous triangularization: expand its matrix
as $R(t)=\sum_\mu t^\mu A_\mu$.  The equations
$R(st)=R(s)R(t)$ and $R(1)=1$ give
$A_\mu A_\eta=0$ for $\mu\ne\eta$, $A_\mu^2=A_\mu$, and
$\sum A_\mu=1$.  Choose a maximal-height weight as in HP2.
Restriction to $1+tE_{ij}$ is polynomial in $t$, since its determinant
is $1$; conjugation by the diagonal torus shows that its coefficient
of $t^a$ raises that weight by $a(e_i-e_j)$.  Hence a maximal-height
weight vector $v$ is fixed by the upper unipotents and transforms
under the diagonal by some integral weight tuple $\lambda$.

Suppose now that the representation $R$ on $W$ is irreducible.
Its orbit of $v$ spans $W$.  For sufficiently large nonnegative
integers $m_\tau$, every matrix coefficient
\[
 f_\phi((X_\tau))=
       \prod_\tau\det(X_\tau)^{m_\tau}\,
                  \phi(R((X_\tau))v),\qquad \phi\in W^*,
 \tag{HP.14}
\]
is polynomial.  These polynomials have right-upper invariance and
column weights $\mu_{\tau,i}=\lambda_{\tau,i}+m_\tau$.
They form a nonzero invariant subspace of the tensor polynomial
space for these column degrees; the map $\phi\mapsto f_\phi$ is
injective because the orbit of $v$ spans $W$.
The column weights are nonnegative because the functions are
polynomial.  HP1--HP2, also in their stated pre-dominance version,
show that $\mu_\tau$ is dominant and that this entire polynomial
space is irreducible.  Thus the coefficient map is onto it.

Its action is exactly
\[
 f_\phi(g^tX)=\prod_\tau\det(g_\tau)^{m_\tau}
                                f_{\phi\circ R(g^t)}(X).
 \tag{HP.15}
\]
Let $\iota(g)=g^{-t}$.  Equation (HP.15) identifies
$(R^*\circ\iota)\otimes\prod\det^{m_\tau}$ with this polynomial
space.  Multiplication by
$\prod\det(X_\tau)^{\mu_{\tau,n}}$ is a nonzero equivariant map
from the canonical tensor module with column degrees
$\mu_{\tau,i}-\mu_{\tau,n}$, with the corresponding determinant
twist, onto that same space.  It is an isomorphism by HP2.
After canceling $\prod\det^{m_\tau}$ the resulting canonical weight
is precisely $\lambda$.  Finally (HP.10) gives the explicit
isomorphism
\[
      V_\lambda\simeq V_\lambda^*\circ\iota,
             \qquad x\longmapsto B(x,-).
 \tag{HP.16}
\]
The operation $R\mapsto R^*\circ\iota$ is involutive, so
$R\simeq V_\lambda$.  The fixed line (HP.7) proves uniqueness of
$\lambda$.  This establishes the required complex classification
and dominance by a matrix-coefficient argument.

The distinction between the two coefficient conventions matters:
\[
\begin{array}{c|c|c}
\text{source representation}&\text{action on rational }\gamma
       &\text{representation whose highest weight is }\lambda\\ \hline
\texttt{tenRepF}(\sigma,\lambda)&\rho_K(\sigma\gamma)
       &\texttt{tenRepF}\ \text{itself}\\
\texttt{wtRepF}(\sigma,\lambda)&\texttt{tenRepF}(\sigma,\lambda)^*
       &\texttt{wtRepF}^{*}\simeq\texttt{tenRepF}\\
\texttt{rhoVr}/\texttt{mtRep}(\sigma,\lambda)
       &\rho_K((\sigma\gamma)^{-t})
       &\texttt{mtRep}^{*}\simeq\texttt{tenRepF}.
\end{array}                                                   \tag{HP.17}
\]
The last equivalence uses the bilinear form (HP.10), and the middle
one is the natural double-dual evaluation map.  For unitary $\gamma$,
$(\sigma\gamma)^{-t}=\sigma c(\gamma)$, so the last line is also
\texttt{tenRep} at the conjugate CM type.  Taking a single dual at
an unchanged embedding would give the usual reversed negative
weight; (HP.17) explains exactly why the source's \emph{dual} weight
is the displayed $\lambda$, with no reversed-negative substitution.

This gives both existence and uniqueness for the literal
\texttt{HWAtF} predicate, including its polynomial-presentation
quantifier.  For existence use the matrices just constructed and the
coordinates of $h$.  Evaluating them at \texttt{evPtF(c,sigma,B)}
puts $B$ in the selected factor, $B^{-t}$ in its conjugate, and the
identity in the others; (HP.5) gives the required Borel eigenvalue.
For uniqueness, any competing polynomial presentation agrees with
the actual one on rational unitary points.  HP3 extends this agreement
to the split group, so it agrees at those exact evaluation points.
A nonzero vector satisfying its Borel equations in one factor may
have arbitrary components in other factors.  Apply a linear functional
to the other tensor factors that does not kill it.  The result is a
nonzero vector in that one polynomial factor, fixed by $U_+$ and of
the proposed diagonal weight.  HP1 and (HP.5) force that weight to be
$\lambda_\tau$.  Distinct Laurent characters are different over an
infinite field.  Thus no change of basis or alternative presentation
can manufacture a different highest weight at the same embedding.

\paragraph{HP5: coefficient descent and the closed CellWt branch.}
\label{claim:HP5}
Let $e:K\hookrightarrow\mathbb C$ and fix a complete CM type
$\sigma:E\hookrightarrow K$ as in AR13.  Let $W$ be an irreducible
finite-dimensional $K$-representation of the rational unitary group,
with the polynomial presentation obtained from its dual's
\texttt{HWAtF} data.  Complexify and choose a nonzero invariant subspace
of minimal dimension.  Apply the classification just proved to its
dual, then dualize back.  Using (HP.16)--(HP.17), this produces
a nonzero complex intertwiner
from a canonical inverse-transpose coefficient module
$V_{\lambda,\mathbb C}$ into $W_\mathbb C$.

The canonical module and its matrices are defined over $K$, by
(HP.3) and the original $\sigma$.  To descend that intertwiner,
choose a $K$-linear functional $\ell:\mathbb C\to K$ not vanishing
on one of its nonzero entries and apply $\ell$ entry by entry.
Every intertwining equation has coefficients in $K$, so the result
is a nonzero $K$-intertwiner.  This argument respects all rational
group elements simultaneously; it is not an unjustified intersection
of infinitely many kernels with scalar extension.
The canonical module is irreducible by HP3, so the descended map
is injective; $W$ is irreducible, so it is surjective.  Hence
\[
       W\simeq_K V_{\lambda,K}
       \quad\text{with its original inverse-transpose action}.
 \tag{HP.18}
\]
The prescribed dual highest weights transfer through this map;
HP4's uniqueness identifies them with this same $\lambda$.
This supplies both representation inputs retained in AR13.
It does not assume that an arbitrary irreducible $K$-representation
has irreducible complexification; in this case that conclusion
follows from (HP.18).

AR14 now completes its original \texttt{CellWt} comparison without
an extra highest-weight premise.  The source's infinitesimal-character
entries $\lambda_i+(n-1)/2-(i-1)$ are strictly decreasing, so equality
of their multisets identifies the dominant weights term by term.
The original \texttt{SideRead} field identifies the coefficient
embeddings, and the canonical polynomial bases commute with $e$.
Conjugating by the actual cell basis-change matrices gives equality
of the representation matrices in those specified bases, exactly
as in AR14.  All other AR14 hypotheses remain: its supplied actual
\texttt{MtReal}, actual \texttt{CharSide} infinitesimal parameters,
coefficient-field maps, and cell base-change witnesses.  The closed
branch is the representation classification and descent used within
that comparison, not the existence of those geometric data.

\paragraph{HP6: the full original-field weight supplier and the new frontier.}
\label{claim:HP6}
The source's \texttt{wt\_weights} theorem is now reconstructed with
no additional representation-theoretic input.  Retain its literal
hypotheses: \texttt{GlobalArtinReciprocity}, $D$, its frame, $j_E$,
\texttt{SideA}, star stability of $D.\mathrm{LL0}$, integral
\texttt{eExp}, and the subsequently quantified original
$\texttt{TwR}/\texttt{RingSpec}$.  The half-integrality needed here is already forced by these records:
$\texttt{cf.isEta.inf}$ says that the framed character has infinite
parameter $1/2$; $\texttt{A.chr.etaT}$ identifies it with the actual
$\eta_T$, and $\texttt{A.chr.infG}$ transports that parameter to
$D.naut$.  Thus a singleton $\{sh\}$ in the antecedent has
$sh=1/2$ (so the source's integer in $sh=z+1/2$ can be $z=0$).
The row-gap and parity calculation, already proved in MR19, gives dominant
$\lambda_j(\tau)$ with
\[
 \operatorname{infOf}(\lambda_j(\tau))
       =\operatorname{paramV}(D.R,D.hc_\tau,j-1,sh)
 \quad\text{when }\operatorname{infG}(\eta_T,\tau)=\{sh\}.
 \tag{HP.19}
\]
The pointwise $sh$ is uniquely determined whenever it exists; where
none exists the implication is empty and the zero weight supplies
the padding.  Use zero weights also outside $1\le j\le m$.

For each original coefficient ring $O_i$, put
\[
 K_i=\operatorname{Frac}(O_i),\qquad
 W_{i,j,b}=\mathrm{VV}_{K_i}(\lambda_j)^*,\qquad
 \rho_{i,j,b}=\texttt{wtRepF}(\texttt{kEmbT}_{i,b},\lambda_j).
 \tag{HP.20}
\]
The exact original side-dependent embeddings are distinct and
nonconjugate.  Explicitly, each is the same injective field reading
composed with the host automorphism labelled by the corresponding
CM-type element.  Equality of two such embeddings would give equality
of the labels; equality with a conjugate would put that same label
in both halves of the CM type.  Both are excluded by the actual frame.
Characteristic zero also follows from the coefficient records: the
finite algebra map $\mathbb Z_\ell\to O_{0,i}$ into a domain DVR is
injective.  Otherwise its prime kernel is $(\ell)$, making $O_{0,i}$
a finite-dimensional domain over $\mathbb F_\ell$, hence a field,
contrary to its DVR nonfield clause.  Composing with the stipulated
injection $O_{0,i}\to O_i$ and the fraction-field injection proves
$\operatorname{char}K_i=0$.  No larger coefficient field is selected.  The five required clauses
now have a single closed supply chain:
\begin{enumerate}
\item HP3 makes the tensor representation irreducible.  Its finite
      dual is irreducible, since the annihilator of a proper nonzero
      dual-stable subspace would be a proper nonzero original-stable
      subspace.  This gives \texttt{rhoW.IsIrreducible}.
\item HP3's explicit polynomial matrices give
      \texttt{IsAlgWeightF} on that very $W_{i,j,b}$.
\item Equation (HP.3), the natural dual-basis map, and injectivity of
      the original coefficient extension identify every
      \texttt{rhoKT} with the same construction over
      $K_{i,k}=\operatorname{Frac}(O_{i,k})$.
      HP3 therefore gives all its irreducibility clauses.
\item The natural double dual and HP4 supply \texttt{HWAtF} for
      $\rho_{i,j,b}^{*}$ at every prescribed embedding, with weights
      $\lambda_j(\tau)$.
\item HP4 forces every weight admitted by that literal predicate to
      equal $\lambda_j(\tau)$; (HP.19) therefore gives the final
      universally quantified infinitesimal-character identity.
\end{enumerate}
There is no representation-dependent extension to put into the lower
bound: one may take $c_2(\ell,j)=\bot$, exactly as in the source.
This closes the original-field, both-side, all-row, all-base-change
\texttt{WtData} supplier, and removes the retained highest-weight
premise from MR19 and from MR5's actual coefficient representation.

The closure is noncircular.  HP1--HP4 use only polynomial linear
algebra, Gaussian elimination, the explicitly proved factorial
pairing, and the Cayley parametrization of the given rational
unitary group.  HP5 then uses the already proved coefficient-extraction
descent and the supplied geometric comparisons.  HP6 uses the already
proved row-weight arithmetic only after the representation theorem is
available.  No \texttt{Core}, theta nonvanishing, automorphic Galois
representation, prime-window estimate, RH or GRH is used in proving
the representation leaf.

For the surrounding construction the frontier has genuinely changed:
\texttt{HwFixedEx}, \texttt{HwFixedLine}, canonical tensor
irreducibility/classification, the conjugate-CM-type dual highest
weight, and weight base-change irreducibility are no longer open
inputs.  \texttt{St/ColPin/R3.Read}, the original tuple and lattice
ties, the good-row spectrum, the supplied \texttt{SideA}, and the
compatible nonzero theta/toric/local installation remain required
where MR/AR use them.  This section supplies none of those by an
unrelated representation of the same dimension.
